\section{Viabilidade e coerência}
\label{sec:coerencia}

Esta seção verifica se o plano é coerente com o restante do projeto e viável com os recursos disponíveis. A Tabela~\ref{tab:coerencia} apresenta as verificações realizadas.

\begin{table}[htb]
\centering
\caption{Verificações de viabilidade e coerência}
\label{tab:coerencia}
\fontsizetable
\begin{tabularx}{\textwidth}{L{3.2cm} Y L{2.2cm}}
\toprule
\textbf{Verificação} & \textbf{Conclusão} & \textbf{Situação}\\
\midrule
Estágio do negócio & As estratégias priorizam validação e densidade local, adequadas à fase; a expansão geográfica é condicionada a gatilhos. & Coerente\\
Modelo de negócio e preços & Os preços (PetCuida+, Plano Parceiro, comissão) e a gratuidade do núcleo seguem a Atividade IX; nenhuma estratégia exige alterar a proposta de valor. & Coerente\\
Canais & A aquisição usa os canais prioritários da Atividade VIII. & Coerente\\
Governança e equipe & Cada ação tem responsável conforme a Atividade XI; o número de frentes simultâneas é limitado a duas estratégias. & Coerente, com atenção à carga\\
Economia unitária & A contribuição de R\$~0,55 por usuário não sustenta escala sem alavancas; o plano trata isso como meta explícita. & Viável sob condições\\
Resultado financeiro & O cenário enxuto é positivo apenas com convênio; o cenário com equipe remunerada exige escala ou financiamento. & Viável sob condições\\
Tecnologia e dados & A infraestrutura remota e a segurança precisam ser concluídas antes de dados reais. & Pré-requisito\\
Aspectos legais & A constituição da empresa e a adequação à LGPD antecedem a coleta de dados de usuários. & Pré-requisito\\
\bottomrule
\end{tabularx}
\fonteautores
\end{table}

\subsection{Condições de continuidade e de revisão}

O plano é viável se três condições se confirmarem no piloto: (i) pelo menos um convênio ou patrocínio é contratado; (ii) o custo de aquisição observado não ultrapassa o teto calculado; e (iii) o fundo de liquidação é reduzido ou limitado, de modo que a margem por atendimento deixe de ser negativa. Se, após seis meses de piloto, o CAC observado exceder o dobro do teto, ou se nenhuma forma de reduzir o fundo se mostrar possível, o grupo deve revisar o modelo antes de ampliar o investimento, por exemplo limitando o uso de créditos, revendo a comissão (faixa de 8\% a 12\% prevista na Atividade IX) ou alterando a cobertura dos créditos solidários.

\subsection{Limitações}

O plano tem limitações que devem ser explicitadas. Primeiro, os valores de custo, preço e volume são estimativas do grupo e não dados observados; a simulação apenas propaga essas premissas. Segundo, mantêm-se constantes os custos fixos, o que subestima o custo de crescer. Terceiro, o modelo não considera a perda de usuários, que elevaria a aquisição bruta necessária. Quarto, os preços de concorrentes citados na Atividade IX são indicativos e sujeitos a confirmação. Por fim, a disposição a pagar das classes C, D e E permanece não medida. Esses limites não invalidam o plano, mas definem o que o piloto precisa medir primeiro.
